\documentclass[11pt]{article}
\usepackage[T1]{fontenc}
\usepackage{textcomp}
\usepackage[margin=0.75in]{geometry}
\usepackage{amsmath,amssymb,amsthm}
\usepackage{array,booktabs}
\usepackage{hyperref}
\setlength{\parindent}{0pt}
\newtheorem{theorem}{Theorem}
\theoremstyle{definition}
\newtheorem{definition}{Definition}
\title{Flow Proof Artifact: Nat-core}
\author{Generated by \texttt{flow doc proof}}
\date{Flow Proof Book}
\begin{document}
\maketitle
Peano predecessor and successor structural laws.\\[0.5em]
\textbf{Source.} peano — https://en.wikipedia.org/wiki/Peano\_axioms\\[0.5em]
\bigskip
\noindent{\large \textbf{Definition 1.} \textit{Taking predecessor after successor returns the same number} \hfill the natural numbers \cdot predecessor \cdot predecessor undoes successor}\par
\smallskip\noindent\textit{Coordinate: the natural numbers \cdot predecessor \cdot predecessor undoes successor}\par
\smallskip\noindent\textit{Source: peano}\par
\medskip
\noindent\textbf{Goal.} Taking predecessor after successor returns the same number.  pred(succ(n)) = n\par
\noindent$\displaystyle \forall n \in \mathbb{N}\quad pred(\mathrm{succ}(n)) = n$\par
\medskip
\noindent
\begin{tabular}{@{} >{\raggedright\arraybackslash}p{0.06\textwidth} >{\raggedright\arraybackslash}p{0.47\textwidth} >{\raggedleft\arraybackslash}p{0.41\textwidth} @{}}
\textbf{\#} & \textbf{Proof} & \textbf{Mathematics} \\
\midrule
\textbf{1.} & We stipulate predecessor undoes successor for predecessor on the natural numbers — this is a definition, not a derived fact. & \\[0.45em]
\textbf{2.} & This follows directly from the definition: pred(the successor of n) equals n. Hence proven. & $\displaystyle pred(\mathrm{succ}(n)) = n$ \\[0.45em]
\bottomrule
\end{tabular}
\label{thm:Nat-predecessor-predecessor_undoes_successor}
\par\medskip\hrule
\bigskip
\noindent{\large \textbf{Definition 2.} \textit{For a nonzero natural, successor after predecessor recovers the number} \hfill the natural numbers \cdot successor \cdot successor undoes predecessor away from zero}\par
\smallskip\noindent\textit{Coordinate: the natural numbers \cdot successor \cdot successor undoes predecessor away from zero}\par
\smallskip\noindent\textit{Source: peano}\par
\medskip
\noindent\textbf{Goal.} For a nonzero natural, successor after predecessor recovers the number\par
\noindent$\displaystyle \forall n \in \mathbb{N}\quad \mathrm{succ}(pred(n)) = n$\par
\medskip
\noindent
\begin{tabular}{@{} >{\raggedright\arraybackslash}p{0.06\textwidth} >{\raggedright\arraybackslash}p{0.47\textwidth} >{\raggedleft\arraybackslash}p{0.41\textwidth} @{}}
\textbf{\#} & \textbf{Proof} & \textbf{Mathematics} \\
\midrule
\textbf{1.} & We stipulate successor undoes predecessor away from zero for successor on the natural numbers — this is a definition, not a derived fact. & \\[0.45em]
\textbf{2.} & We split into exhaustive cases — the claim must hold in each one. & \\[0.45em]
\textbf{3.} & Case 1 (see step 2): suppose n is zero. & \\[0.45em]
\textbf{4.} & From step 3, this implies the successor of pred(n) equals n in this case. & $\displaystyle \mathrm{succ}(pred(n)) = n$ \\[0.45em]
\textbf{5.} & Case 2 (see step 2): neither disjunct holds. & \\[0.45em]
\textbf{6.} & Let k denote the predecessor of n. & \\[0.45em]
\textbf{7.} & We invoke Definition 1 (predecessor undoes successor, for predecessor on the natural numbers, instantiated for k). & \\[0.45em]
\textbf{8.} & From step 5, step 6, and Definition 1, this implies the successor of pred(n) equals n. Together with the other cases (step 3 and step 5), the goal is discharged. Hence proven. & $\displaystyle \mathrm{succ}(pred(n)) = n$ \\[0.45em]
\bottomrule
\end{tabular}
\label{thm:Nat-successor-successor_undoes_predecessor_away_from_zero}
\par\medskip\hrule
\bigskip
\noindent{\large \textbf{Derived fact 3.} \textit{If two successors match, the originals match} \hfill the natural numbers \cdot successor \cdot successor is injective}\par
\smallskip\noindent\textit{Coordinate: the natural numbers \cdot successor \cdot successor is injective}\par
\smallskip\noindent\textit{Source: peano — https://en.wikipedia.org/wiki/Injective\_function}\par
\smallskip\noindent\textit{Built on: \hyperref[thm:Nat-predecessor-predecessor_undoes_successor]{Definition 1}}\par
\medskip
\noindent\textbf{Goal.} If two successors match, the originals match.  succ(m) = succ(n) → m = n\par
\noindent$\displaystyle \forall m \in \mathbb{N} \forall n \in \mathbb{N}\quad m = n$\par
\medskip
\noindent
\begin{tabular}{@{} >{\raggedright\arraybackslash}p{0.06\textwidth} >{\raggedright\arraybackslash}p{0.47\textwidth} >{\raggedleft\arraybackslash}p{0.41\textwidth} @{}}
\textbf{\#} & \textbf{Proof} & \textbf{Mathematics} \\
\midrule
\textbf{1.} & We prove that successor is injective for successor on the natural numbers. & \\[0.45em]
\textbf{2.} & We split into exhaustive cases — the claim must hold in each one. & \\[0.45em]
\textbf{3.} & Case 1 (see step 2): suppose succ(m)  equals  succ(n). & \\[0.45em]
\textbf{4.} & We invoke Definition 1 (predecessor undoes successor, for predecessor on the natural numbers, instantiated for m). & \\[0.45em]
\textbf{5.} & We invoke Definition 1 (predecessor undoes successor, for predecessor on the natural numbers, instantiated for n). & \\[0.45em]
\textbf{6.} & From step 3, Definition 1, and Definition 1, this implies m equals n. Together with the other cases (step 3), the goal is discharged. Hence proven. & $\displaystyle m = n$ \\[0.45em]
\bottomrule
\end{tabular}
\label{thm:Nat-successor-successor_is_injective}
\par\medskip\hrule
\bigskip
\noindent{\large \textbf{Derived fact 4.} \textit{Zero is never the successor of any natural number} \hfill the natural numbers \cdot zero \cdot zero is not a successor}\par
\smallskip\noindent\textit{Coordinate: the natural numbers \cdot zero \cdot zero is not a successor}\par
\smallskip\noindent\textit{Source: peano}\par
\smallskip\noindent\textit{Built on: \hyperref[thm:Nat-predecessor-predecessor_undoes_successor]{Definition 1}}\par
\medskip
\noindent\textbf{Goal.} Zero is never the successor of any natural number\par
\noindent$\displaystyle \forall n \in \mathbb{N}\quad pred(0) = n$\par
\medskip
\noindent
\begin{tabular}{@{} >{\raggedright\arraybackslash}p{0.06\textwidth} >{\raggedright\arraybackslash}p{0.47\textwidth} >{\raggedleft\arraybackslash}p{0.41\textwidth} @{}}
\textbf{\#} & \textbf{Proof} & \textbf{Mathematics} \\
\midrule
\textbf{1.} & We prove that zero is not a successor for zero on the natural numbers. & \\[0.45em]
\textbf{2.} & We split into exhaustive cases — the claim must hold in each one. & \\[0.45em]
\textbf{3.} & Case 1 (see step 2): suppose 0  equals  succ(n). & \\[0.45em]
\textbf{4.} & We invoke Definition 1 (predecessor undoes successor, for predecessor on the natural numbers, instantiated for n). & \\[0.45em]
\textbf{5.} & From step 3 and Definition 1, this implies pred(0) equals n. Together with the other cases (step 3), the goal is discharged. Hence proven. & $\displaystyle pred(0) = n$ \\[0.45em]
\bottomrule
\end{tabular}
\label{thm:Nat-zero-zero_is_not_a_successor}
\par\medskip\hrule
\bigskip
\noindent{\large \textbf{Derived fact 5.} \textit{Successor never produces zero} \hfill the natural numbers \cdot successor \cdot successor never yields zero}\par
\smallskip\noindent\textit{Coordinate: the natural numbers \cdot successor \cdot successor never yields zero}\par
\smallskip\noindent\textit{Source: peano}\par
\smallskip\noindent\textit{Built on: \hyperref[thm:Nat-zero-zero_is_not_a_successor]{Derived fact 4}}\par
\medskip
\noindent\textbf{Goal.} Successor never produces zero\par
\noindent$\displaystyle \forall n \in \mathbb{N}\quad \mathrm{succ}(n) != 0$\par
\medskip
\noindent
\begin{tabular}{@{} >{\raggedright\arraybackslash}p{0.06\textwidth} >{\raggedright\arraybackslash}p{0.47\textwidth} >{\raggedleft\arraybackslash}p{0.41\textwidth} @{}}
\textbf{\#} & \textbf{Proof} & \textbf{Mathematics} \\
\midrule
\textbf{1.} & We prove that successor never yields zero for successor on the natural numbers. & \\[0.45em]
\textbf{2.} & We split into exhaustive cases — the claim must hold in each one. & \\[0.45em]
\textbf{3.} & Case 1 (see step 2): suppose succ(n)  equals  0. & \\[0.45em]
\textbf{4.} & We invoke Derived fact 4 (zero is not a successor, for zero on the natural numbers, instantiated for n). & \\[0.45em]
\textbf{5.} & From step 3 and Derived fact 4, this implies the successor of n ! equals 0 in this case. & $\displaystyle \mathrm{succ}(n) != 0$ \\[0.45em]
\textbf{6.} & Case 2 (see step 2): neither disjunct holds. & \\[0.45em]
\textbf{7.} & From step 6, this implies the successor of n ! equals 0. Together with the other cases (step 3 and step 6), the goal is discharged. Hence proven. & $\displaystyle \mathrm{succ}(n) != 0$ \\[0.45em]
\bottomrule
\end{tabular}
\label{thm:Nat-successor-successor_never_yields_zero}
\par\medskip\hrule
\bigskip
\noindent{\large \textbf{Derived fact 6.} \textit{Every natural is either zero or the successor of some natural} \hfill the natural numbers \cdot cases \cdot every number is zero or a successor}\par
\smallskip\noindent\textit{Coordinate: the natural numbers \cdot cases \cdot every number is zero or a successor}\par
\smallskip\noindent\textit{Source: peano}\par
\smallskip\noindent\textit{Built on: \hyperref[thm:Nat-predecessor-predecessor_undoes_successor]{Definition 1}}\par
\medskip
\noindent\textbf{Goal.} Every natural is either zero or the successor of some natural\par
\noindent$\displaystyle \forall n \in \mathbb{N}\quad n = \mathrm{succ}(k)$\par
\medskip
\noindent
\begin{tabular}{@{} >{\raggedright\arraybackslash}p{0.06\textwidth} >{\raggedright\arraybackslash}p{0.47\textwidth} >{\raggedleft\arraybackslash}p{0.41\textwidth} @{}}
\textbf{\#} & \textbf{Proof} & \textbf{Mathematics} \\
\midrule
\textbf{1.} & We prove that every number is zero or a successor for cases on the natural numbers. & \\[0.45em]
\textbf{2.} & We split into exhaustive cases — the claim must hold in each one. & \\[0.45em]
\textbf{3.} & Case 1 (see step 2): suppose n is zero. & \\[0.45em]
\textbf{4.} & From step 3, this implies n equals 0 in this case. & $\displaystyle n = 0$ \\[0.45em]
\textbf{5.} & Case 2 (see step 2): neither disjunct holds. & \\[0.45em]
\textbf{6.} & Let k denote the predecessor of n. & \\[0.45em]
\textbf{7.} & We invoke Definition 1 (predecessor undoes successor, for predecessor on the natural numbers, instantiated for k). & \\[0.45em]
\textbf{8.} & From step 5, step 6, and Definition 1, this implies n equals the successor of k. Together with the other cases (step 3 and step 5), the goal is discharged. Hence proven. & $\displaystyle n = \mathrm{succ}(k)$ \\[0.45em]
\bottomrule
\end{tabular}
\label{thm:Nat-cases-every_number_is_zero_or_a_successor}
\par\medskip\hrule
\end{document}
